\documentclass[10pt, conference]{FMFP2023}

% ---------- FMFP Required Packages ----------
\usepackage[top=1.50cm, bottom=1.50cm, left=1.884cm, right=1.884cm,
            includehead, includefoot, heightrounded]{geometry}
\usepackage{fancyhdr}
\usepackage{graphicx}
\usepackage[textfont=bf, font=normalsize]{caption}
\usepackage{floatrow}
\usepackage{amssymb}
\usepackage{amsfonts}
\usepackage{amsmath}
\floatstyle{plaintop}
\restylefloat{table}
\usepackage{subcaption}   % replaces deprecated subfigure package
\usepackage{booktabs}     % professional table rules
\floatsetup{font=normalsize}

% ---------- User Specific Packages ----------
\usepackage{siunitx}
\usepackage{physics}
\usepackage{bm}
\usepackage{hyperref}
\usepackage{enumitem}
\usepackage[ruled,vlined,linesnumbered]{algorithm2e}
\usepackage{nicefrac}

\hypersetup{
    colorlinks=true,
    linkcolor=blue,
    citecolor=blue,
    urlcolor=blue
}

% ---------- FMFP Header Configuration ----------
\rhead{%
    \sffamily\fontsize{9}{11}\selectfont
    \vspace{-0.5cm}
    \textbf{Proceedings of the 10th International and 50th National
    Conference on Fluid Mechanics and Fluid Power (FMFP) \\
    December 20--22, 2023, IIT Jodhpur, Rajasthan, India}
    \\ \vspace{0.15cm}
    \large{\textbf{FMFP2023-XX-YYY}}}

\renewcommand{\headrulewidth}{0pt}

\begin{document}

% ---------- Title & Authors ----------
\title{\Large\textbf{GPU-Accelerated Implementation of the
Savage--Hutter Model with Application to 1D and 2D Inclined Plane
Granular Flows}}

\author{\textbf{Samyak Mayank}\textsuperscript{1}\\
\textsuperscript{1}\small{School of Mechanical and Materials
Engineering, Indian Institute of Technology Mandi,
Himachal Pradesh 175075, India}}

\maketitle
\thispagestyle{fancy}
\pagestyle{plain}

% ---------- Abstract ----------
\noindent\textbf{ABSTRACT}

\noindent
Granular flows such as rockfalls, landslides, and snow avalanches pose
serious hazards in mountainous regions and are characterized by complex
rheology, intermediate between solid-like and fluid-like behavior.
This paper models the motion of a finite mass of cohesionless granular
material released from rest on a rough inclined plane using the
depth-averaged Savage--Hutter equations.
Numerical simulations are validated against analytical similarity
solutions (parabolic cap profile) and then accelerated using CUDA-based
GPU parallelization.
A performance comparison between Python, C++ (CPU), and CUDA (GPU)
implementations highlights the computational advantages of GPU computing
for large-scale granular flow simulations.
The framework is further extended to a full two-dimensional
finite-volume Rusanov solver, rigorously validated against the TITAN2D
reference code using Intersection-over-Union (IoU), RMSE, and
Lagrangian centroid tracking metrics.

\smallskip
\noindent\textbf{Keywords:} Granular Avalanches, Savage--Hutter Model,
GPU Acceleration, CUDA, Finite Difference Method, Finite Volume Method

% ============================================================
% FIX: removed \textbf{} wrappers from all \section and \subsection
%      commands — the class already bolds section headings, so wrapping
%      in \textbf produces double-bold artefacts.
% ============================================================
\section{Introduction}

Granular flows such as rockfalls, landslides, and snow avalanches
constitute major geological hazards in mountainous regions.
Unlike simple Newtonian fluids, granular materials exhibit dual
behavior: they can flow like fluids when mobilized, yet they also
possess solid-like frictional strength that allows them to remain at
rest on steep slopes.
This combination of frictional resistance, dilatancy, and collisional
interactions makes granular avalanche modeling non-trivial.

In order to investigate the fundamentals of avalanche motion, it is
useful to focus on an idealized configuration: a finite mass of
cohesionless granular material released from rest on a rough inclined
plane.
After release, the mass accelerates downslope under gravity, spreads
due to internal pressure gradients, and eventually comes to rest as
basal friction balances the driving component of gravity.
The objective of this paper is threefold:
\begin{itemize}[nosep]
    \item To model the avalanche using depth-averaged continuum
          equations based on the Savage--Hutter framework.
    \item To implement a numerical finite-difference scheme and
          validate the simulation against the analytical parabolic cap
          similarity solution.
    \item To introduce GPU acceleration using CUDA and demonstrate the
          performance advantage of GPU-based implementation over Python
          and C++ CPU-based codes.
\end{itemize}

% Literature Review section removed to conserve page budget.

% ============================================================
\section{Problem Definition}
% ============================================================

\subsection{Physical Setup}

We consider a finite mass of cohesionless granular material initially
at rest on a rigid, planar bed inclined at an angle $\zeta$ to the
horizontal.
The granular body has an initial thickness distribution $h(x,0)$,
nonzero over a finite interval, and zero elsewhere.
When the material is released, gravity accelerates the mass downslope,
internal pressure gradients drive spreading, and basal friction
decelerates and eventually stops the flow.
The free surface and the positions of the front and rear margins evolve
in time, making this a moving-boundary problem.

\subsection{Modeling Assumptions}

\paragraph{Incompressibility.}
The bulk density $\rho$ is taken to be constant and changes in void
ratio during flow are neglected, so mass conservation simplifies to
depth conservation.

\paragraph{Shallow Flow.}
The flow is assumed to be long and shallow, with a small aspect ratio
\[
  \epsilon = \frac{H}{L} \ll 1,
\]
where $H$ is a characteristic thickness and $L$ a characteristic
length along the slope.
This justifies using a hydrostatic normal stress approximation and
performing depth-averaging.

\paragraph{Depth Averaging.}
Velocity and stress variations in the vertical direction are
represented through their depth-averaged counterparts.
The primary kinematic variables are therefore $h(x,t)$ (flow thickness)
and $u(x,t)$ (depth-averaged downslope velocity).

\paragraph{Frictional Resistance.}
Basal friction is modeled by a Coulomb sliding law with constant basal
friction angle $\delta$: $\tau_b = \sigma_n\tan\delta$, where $\tau_b$
is the basal shear stress and $\sigma_n$ the normal stress at the bed.
Internal friction is characterized by an angle $\phi$, which enters
the earth-pressure coefficient.

% ============================================================
\section{Mathematical Formulation}
% ============================================================

Under the above assumptions and after non-dimensionalization, the
depth-averaged governing equations for flow over a planar bed take the
following form.

\subsection{Conservation of Mass}

\begin{equation}
    \frac{\partial h}{\partial t}
    + \frac{\partial}{\partial x}(h\,u) = 0,
    \label{eq:mass}
\end{equation}
which expresses conservation of mass (or thickness) for the deforming
granular body.

\subsection{Conservation of Momentum}

\begin{equation}
    \frac{\partial u}{\partial t}
    + u\,\frac{\partial u}{\partial x}
    = \sin\zeta
    - \tan\delta\cos\zeta
    - \beta\,\frac{\partial h}{\partial x},
    \label{eq:momentum}
\end{equation}
where $\sin\zeta$ is the driving downslope gravity component,
$\tan\delta\cos\zeta$ is the resisting basal Coulomb friction (for
$u>0$), and $\beta\,\partial h/\partial x$ is an effective
pressure-gradient term.
The coefficient
\begin{equation}
    \beta = \epsilon\,k_{\text{act}}\cos\zeta,
\end{equation}
where $\epsilon$ is the aspect ratio and $k_{\text{act}}$ is the
active earth-pressure coefficient.

\subsection{Earth-Pressure Coefficient}

The earth-pressure coefficient $k_{\text{act}}$ is derived from the
Mohr--Coulomb yield criterion under an active stress state (material
effectively extending downslope).
It depends primarily on the internal friction angle $\phi$ and slope
angle $\zeta$, and is treated as a constant for a given material and
slope.

\subsection{Boundary Conditions}

Because the granular mass is finite, the thickness vanishes at the
front and rear margins:
\begin{equation}
    h(x_F(t),t) = 0, \qquad h(x_R(t),t) = 0,
\end{equation}
with margins moving at the local flow velocity
$dx_{F,R}/dt = u(x_{F,R},t)$.
In numerical simulations over a fixed domain $[L_{\text{left}},
L_{\text{right}}]$, we impose $h=0$, $u=0$ at both far-field
boundaries.

% ============================================================
\section{Solution Techniques}
% ============================================================

\subsection{Analytical Similarity Solutions}

For special classes of initial conditions and under symmetry
assumptions, the governing equations admit analytical self-similar
solutions, which serve as benchmarks for validating the numerical
implementation.

\subsubsection{Parabolic Cap Solution}

A particularly important similarity solution is the
\emph{parabolic cap}.
In a suitably scaled comoving coordinate
$\eta = (x - x_c(t))/g(t)$, where $x_c(t)$ is the center of mass
and $g(t)$ is the half-width of the pile, the thickness profile takes
the form
\begin{equation}
    h(\eta,t) = \frac{K}{g(t)}\bigl(1-\eta^2\bigr),
    \label{eq:parabolic-cap}
\end{equation}
where $K$ is a constant related to the conserved total mass.
The depth-averaged velocity varies approximately linearly with space,
$\tilde{u}(\eta,t)\propto\eta$, with maximum magnitude near the center
and zero at the margins.
The half-width grows asymptotically as $g(t)\sim t$ as $t\to\infty$,
while the maximum height decays due to spreading.

\subsection{Numerical Methods}

For general initial conditions, the system~\eqref{eq:mass}--\eqref{eq:momentum}
must be solved numerically.
The computational domain $[L_{\text{left}}, L_{\text{right}}]$ is
discretized into $N_x$ grid points $x_i$, and time is advanced using
an explicit forward Euler scheme:
\begin{equation}
    q_i^{n+1} = q_i^n + \Delta t\,\mathcal{R}_i(q^n),
\end{equation}
where $q=(h,u)$ and $\mathcal{R}$ denotes the discrete right-hand
side.
Stability requires a CFL condition relating $\Delta t$ to mesh size
and characteristic wave speeds.
The convective term $u\,\partial u/\partial x$ is approximated using
an upwind finite-difference scheme for stability, while
$\partial h/\partial x$ uses centered differences to avoid artificial
bias in the spreading direction.
At the far-field boundaries, $h=0$ and $u=0$ are imposed directly.

% ============================================================
\section{Simulation Scenarios and Validation Strategy}
% ============================================================

Three simulation cases are used to systematically verify the numerical
solver, each isolating a specific part of the coupled
mass--momentum system.

\subsection{Case 1: Numerical Thickness, Analytical Velocity}

The analytical similarity velocity profile is prescribed while the
mass conservation equation is solved numerically to obtain $h(x,t)$.
This isolates the spatial discretization of the mass equation and
confirms correct spreading dynamics and mass preservation.
The numerical thickness profiles relax into a parabolic shape,
with front and rear margins advancing outward at nearly constant rates
for large $t$, matching the asymptotic trend $g(t)\sim t$.
Small deviations near the boundaries arise from numerical diffusion
inherent to first-order upwinding.

\subsection{Case 2: Numerical Velocity, Analytical Thickness}

The analytical thickness profile is imposed while the momentum
equation is solved numerically for $u(x,t)$.
This evaluates the accuracy of the momentum solver independently,
testing advection stability and internal-pressure effects.
The numerically computed velocity distribution maintains nearly linear
spatial variation $u(x,t)\propto\eta$, as predicted by the similarity
solution, with minor phase lag that decreases under grid refinement.

\subsection{Case 3: Fully Coupled Numerical Solution}

Both mass and momentum equations are solved simultaneously using the
finite difference scheme (Fig.~\ref{fig:case3_coupled}), representing
the full physical problem and assessing coupling stability and
convergence to similarity form.

\begin{figure*}[t!]
    \centering
    \includegraphics[width=0.95\textwidth]{coupled.pdf}
    \caption{Case 3: Fully coupled numerical profiles of $h(x,t)$
    (left) and $u(x,t)$ (right) compared with the analytical
    Savage--Hutter parabolic cap similarity solution
    ($\zeta=32^\circ$, $\delta=22^\circ$, $\beta=0.204$).}
    \label{fig:case3_coupled}
\end{figure*}

This staged validation structure gives strong evidence that all solver
components are correct before GPU acceleration is introduced.

% ============================================================
\section{Results and Discussion}
% ============================================================

The results demonstrate that the finite difference solver successfully
reproduces the spreading and flattening behavior predicted by the
Savage--Hutter model.

\subsection{Case 1: Mass Conservation Verification}

The numerical thickness profiles quickly relax into a parabolic shape.
The front and rear margins advance outward at nearly constant rates
for large $t$, matching the asymptotic spreading trend $g(t)\sim t$.
Small deviations near boundaries are due to numerical diffusion
inherent to first-order upwinding.

\subsection{Case 2: Momentum Solver Verification}

The numerically computed velocity distribution maintains nearly linear
spatial variation $u(x,t)\propto\eta$, as predicted by the similarity
solution.
Minor phase lag between numerical and analytical profiles decreases
with grid refinement, confirming correct implementation of
internal-pressure and advective transport terms.

\subsection{Case 3: Fully Coupled Dynamics}

In the coupled case, the solution converges naturally toward the
parabolic cap profile even when initial conditions are not perfectly
parabolic, providing numerical confirmation of the theoretical
\emph{shape stability} of the solution.
The spreading rate and profile alignment with analytical similarity
solutions indicate strong accuracy of the coupled solver.

% ============================================================
\section{GPU Programming and Performance}
% ============================================================

\subsection{Motivation for GPU Acceleration}

As spatial resolution and simulation time increase, the computational
cost of solving the Savage--Hutter equations on a CPU becomes
prohibitive.
GPUs are designed for massively parallel computations and can execute
thousands of lightweight threads simultaneously, making them
well-suited for updating large arrays of $h$ and $u$ in parallel,
applying stencil-based finite-difference operations, and running many
time steps efficiently.

\subsection{Methodology of GPU Acceleration}

\subsubsection{Conceptual Overview}

In a standard CPU (serial) implementation, a single processing unit
iterates through the grid cell-by-cell, computing inter-cell fluxes
and updating each cell sequentially.
As grid resolution scales to hundreds of thousands of cells, this
sequential processing becomes a profound bottleneck.
GPU acceleration takes a fundamentally different approach: the GPU
deploys one thread per grid cell, and at each time step all threads
simultaneously read their neighbors, compute fluxes, and update their
respective cells concurrently.

\subsubsection{Algorithmic Framework: Serial vs.\ Parallel}

The transition from a sequential CPU model to a parallel CUDA model
involves flattening the nested spatial loops.
The explicit nature of the Rusanov scheme ensures that the update for
cell $(i,j)$ depends only on its neighbors at the current time step
$t^n$, making the algorithm \emph{embarrassingly parallel}.
% FIX: label added to both algorithms for correct cross-referencing.
%      Algorithm 1 previously referenced Eq.~\ref{eq:fv_update} which
%      is defined later in Section 9; the label is retained and the
%      forward reference is now resolved correctly.
\begin{algorithm}[t]
\caption{Serial CPU Execution}
\label{alg:cpu}
\SetAlgoLined
Allocate memory for state vector $\bm{U} = (h,\,hu,\,hv)$\;
\For{$t = 0$ \KwTo $T_{\text{end}}$}{
    Compute maximum wave speed and $\Delta t$\;
    \For{$j = 1$ \KwTo $N_y$}{
        \For{$i = 1$ \KwTo $N_x$}{
            Compute $k(x,y)$ using velocity divergence\;
            Compute interface fluxes $\hat{\bm{F}}_{i+1/2,j}$
            and $\hat{\bm{G}}_{i,j+1/2}$\;
        }
    }
    \For{$j = 1$ \KwTo $N_y$}{
        \For{$i = 1$ \KwTo $N_x$}{
            Update state $\bm{U}^{n+1}_{i,j}$
            via Eq.~\eqref{eq:fv_update}\;
            Enforce dry-bed constraint ($h \ge 0$)\;
        }
    }
}
\end{algorithm}

\begin{algorithm}[t]
\caption{Parallel GPU Execution (Host-Side Coordination)}
\label{alg:gpu}
\SetAlgoLined
Allocate device memory using Struct of Arrays (SoA) layout\;
Copy initial conditions from host to device\;
\For{$t = 0$ \KwTo $T_{\text{end}}$}{
    Launch \texttt{compute\_dt} reduction kernel\;
    Launch \texttt{compute\_k\_grid} kernel
    $\langle\!\langle\!\langle\text{Grid},\,\text{Block}
    \rangle\!\rangle\!\rangle$\;
    Launch \texttt{compute\_fluxes} kernel
    $\langle\!\langle\!\langle\text{Grid},\,\text{Block}
    \rangle\!\rangle\!\rangle$\;
    Launch \texttt{update\_state} kernel
    $\langle\!\langle\!\langle\text{Grid},\,\text{Block}
    \rangle\!\rangle\!\rangle$\;
}
Copy final results from device to host for visualization\;
\end{algorithm}

\subsubsection{Spatial Mapping and Kernel Execution}

The two-dimensional computational domain is partitioned into a uniform
grid of CUDA thread blocks (e.g., $16\times16$ threads per block).
Physics computations are distributed across distinct
\texttt{\_\_global\_\_} kernel functions to enforce global memory
synchronization before advancing to subsequent operations.
Inside each kernel, a thread determines its unique global spatial
indices:
\begin{align}
    i &= \text{blockIdx.x} \times \text{blockDim.x}
         + \text{threadIdx.x}, \\
    j &= \text{blockIdx.y} \times \text{blockDim.y}
         + \text{threadIdx.y}.
\end{align}

\subsubsection{Data Layout and Memory Optimization}

Traditional object-oriented programming groups variables into an
Array of Structures (AoS).
The GPU implementation instead uses a Struct of Arrays (SoA), storing
all flow thicknesses $h$ contiguously in memory, followed separately
by momentum arrays $hu$ and $hv$.
When a CUDA warp (32 synchronized threads) accesses SoA memory, it
performs a single coalesced memory transaction, saturating the
high-bandwidth GDDR memory bus.
Host-to-device transfers (\texttt{cudaMemcpy}) occur only at
initialization and specific checkpoints for data logging, eliminating
PCIe bus latency during the temporal loop.

\subsection{Performance Comparison and Analysis}

The numerical solver was executed for grid resolutions from
$N_x=101$ to $N_x=100{,}001$, measuring execution time for Python
(NumPy-based), C++ (CPU), and CUDA (GPU) implementations.

% FIX: Execution time figure converted to figure* (full-width).
%      Previously a single-column figure with width=1\textwidth, which
%      overflowed into the adjacent column.  Full-width is required to
%      make the many grid-size labels on the x-axis legible.
\begin{figure*}[t!]
    \centering
    \includegraphics[width=0.95\textwidth]{clean_execution_time.pdf}
    \caption{Execution time for Python (green), C++ (orange), and
    CUDA (blue) implementations across different 1D grid sizes $N_x$.
    Python is slowest due to its interpreted nature; the GPU incurs
    overhead for small grids but dominates at large $N_x$.}
    \label{fig:time-vs-size}
\end{figure*}

% FIX: Speedup histogram converted to figure* for readability.
\begin{figure*}[t!]
    \centering
    \includegraphics[width=0.92\textwidth]{speedup_histogram.pdf}
    \caption{Speedup achieved by compiler optimization levels
    \texttt{-O1} (blue), \texttt{-O2} (orange), and \texttt{-O3}
    (green) relative to unoptimized execution (\texttt{-O0}) as a
    function of 2D grid size. Both \texttt{-O2} and \texttt{-O3}
    achieve the highest gains.}
    \label{fig:opt_speedup}
\end{figure*}

% FIX: CUDA vs C++ speedup converted to figure* for readability.
\begin{figure*}[t!]
    \centering
    \includegraphics[width=0.95\textwidth]{clean_speedup_cpp.pdf}
    \caption{Speedup of GPU (CUDA) relative to the optimized C++
    implementation as a function of 1D grid size $N_x$. Values
    greater than 1 indicate GPU outperforming CPU; the crossover
    occurs near $N_x \approx 25{,}601$.}
    \label{fig:speedup_cpp}
\end{figure*}

% FIX: CUDA vs Python speedup converted to figure* for readability.
\begin{figure*}[t!]
    \centering
    \includegraphics[width=0.95\textwidth]{clean_speedup_python.pdf}
    \caption{Speedup of GPU (CUDA) relative to the Python (NumPy)
    implementation as a function of 1D grid size $N_x$.
    The GPU achieves more than $70\times$ speedup at $N_x=1601$.}
    \label{fig:speedup_python}
\end{figure*}

As shown in Fig.~\ref{fig:time-vs-size}, Python performs slowest
across all grid sizes.
The C++ CPU implementation scales far more efficiently and is
significantly faster for small to medium grid sizes
($N_x < 25{,}000$).
The GPU initially exhibits larger runtimes for smaller grids because
kernel launch latency and host--device memory transfers dominate.
However, as $N_x$ increases, the advantage of massive parallelism
becomes evident (Figs.~\ref{fig:speedup_cpp}
and~\ref{fig:speedup_python}).
The GPU achieves more than $70\times$ speedup over Python at
$N_x=1601$, and outperforms C++ for grid sizes exceeding approximately
$N_x=25{,}601$.
Compiler optimization levels for 2D grids are benchmarked in
Fig.~\ref{fig:opt_speedup}, with \texttt{-O2} and \texttt{-O3}
achieving the highest gains over the unoptimized baseline.
These results confirm that CUDA acceleration enables simulations at
resolutions that would be impractical on CPU-based solvers alone.

% ============================================================
% CORRECTED NEW SECTIONS — v2
% Four sections to insert AFTER Section VIII (GPU Programming
% and Performance) and BEFORE the "Extension to Two-Dimensional
% Flow" section.
%
% Fixes applied vs. v1 / vs. compiled PDF:
%  1. Cross-reference "Section~9" → \ref{sec:2d_governing} so it
%     resolves correctly regardless of final section numbering.
%  2. Algorithm counter: old paper already has Alg. 1 (Serial CPU)
%     and Alg. 2 (Parallel GPU) for the 1D solver.  The four new
%     algorithms are therefore labelled alg:cpu_2d (→ Alg. 3) and
%     alg:gpu_2d (→ Alg. 4) — matching what appears in the PDF.
%  3. Roofline citation [?] replaced with a proper \bibitem key
%     \cite{Williams2009}; add the entry below to \thebibliography.
%  4. Figure placeholders (nsight_throughput.pdf, strong_scaling_2d.pdf
%     weak_scaling_2d.pdf, soa_vs_aos.pdf) retained — they already
%     appear as empty boxes in the compiled PDF awaiting real figures.
%  5. Wording tightened to match the FMFP two-column IEEE style used
%     throughout the document (no \paragraph{} inside subsections —
%     use \subsubsection or plain prose instead; FMFP class renders
%     \paragraph as run-in bold which conflicts with subsection style).
%  6. "Section 9" hard-code removed; label \label{sec:2d_governing}
%     must be placed on the "Extension to Two-Dimensional Flow"
%     section heading in the main document.
%  7. Table column alignment corrected: booktabs + @{} flush,
%     matching Tables 1 and 3 in the compiled paper.
%  8. Added \label{sec:gpu2d}, \label{sec:roofline}, \label{sec:scaling},
%     \label{sec:datalayout} to each new section for internal \ref use.
% ============================================================

% ------------------------------------------------------------------
% REQUIRED ADDITION TO \thebibliography IN MAIN DOCUMENT:
%
% \bibitem{Williams2009}
% Williams, S., Waterman, A., \& Patterson, D. (2009).
% \textit{Roofline: an insightful visual performance model for
% multicore architectures}.
% \textit{Communications of the ACM}, 52(4), 65--76.
% ------------------------------------------------------------------

% ============================================================
\section{GPU Parallelisation of the Two-Dimensional Solver}
\label{sec:gpu2d}
% ============================================================

\subsection{Kernel Decomposition and Thread Mapping}

The two-dimensional Rusanov solver described in
Section~\ref{sec:2d_governing} requires the evaluation of
four logically distinct operations per time step:
(i) computation of the dynamic earth-pressure field
$k(x,y,t)$,
(ii) estimation of the global maximum wave speed for
adaptive time-stepping,
(iii) computation of all inter-cell Rusanov fluxes, and
(iv) update of the conserved state vector
$\bm{U}=(h,hu,hv)^\top$.
Each operation is implemented as a separate
\texttt{\_\_global\_\_} CUDA kernel, and the mandatory
\texttt{cudaDeviceSynchronize} barrier between successive
launches enforces the data dependency chain without
requiring explicit software barriers inside any kernel.

\subsubsection{Algorithmic Framework}

The transition from a sequential CPU model to a massively
parallel CUDA model requires restructuring the solver so
that every cell $(i,j)$ maps to a unique GPU thread.
The update for cell $(i,j)$ at time step $n+1$ depends only
on the values of its four face-adjacent neighbours at time
step $n$; this nearest-neighbour stencil pattern makes the
explicit Rusanov scheme \emph{embarrassingly parallel} with
no cyclic data dependence.
Algorithms~\ref{alg:cpu_2d} and~\ref{alg:gpu_2d} summarise
the serial CPU and parallel GPU execution paths,
respectively.

\begin{algorithm}[t]
\caption{Serial 2D CPU Rusanov Execution}
\label{alg:cpu_2d}
\SetAlgoLined
Allocate arrays $h$, $hu$, $hv$, $k_{\text{grid}}$ of
size $N_x \times N_y$\;
Initialise hemispherical cap: $h(x,y,0)$ from
Eq.~\eqref{eq:initial_condition}\;
\For{$t = 0$ \KwTo $T_{\text{end}}$}{
    \For{$j = 1$ \KwTo $N_y - 1$}{
        \For{$i = 1$ \KwTo $N_x - 1$}{
            Compute divergence
            $\nabla\cdot\bar{\bm{u}}_{i,j}$\;
            Assign $k_{i,j} \leftarrow k_{\text{act}}$
            or $k_{\text{pass}}$\;
        }
    }
    Find $a_{\max} = \max_{i,j}(|u_{i,j}|+c_{i,j},\;
    |v_{i,j}|+c_{i,j})$; set $\Delta t$\;
    \For{$j = 0$ \KwTo $N_y$}{
        \For{$i = 0$ \KwTo $N_x$}{
            Compute $\hat{\bm{F}}_{i+\nicefrac{1}{2},j}$
            and $\hat{\bm{G}}_{i,j+\nicefrac{1}{2}}$\;
        }
    }
    \For{$j = 1$ \KwTo $N_y - 1$}{
        \For{$i = 1$ \KwTo $N_x - 1$}{
            Update $\bm{U}^{n+1}_{i,j}$ via
            Eq.~\eqref{eq:fv_update}\;
            Enforce dry-bed constraint ($h \geq 0$)\;
        }
    }
}
\end{algorithm}

\begin{algorithm}[t]
\caption{Parallel 2D GPU (CUDA) Execution}
\label{alg:gpu_2d}
\SetAlgoLined
\textbf{Host:} Allocate SoA device arrays for $h$, $hu$,
$hv$, $k_{\text{grid}}$, flux buffers
$\bm{F}^x_{h,hu,hv}$, $\bm{G}^y_{h,hu,hv}$,
\texttt{wave\_speeds}\;
\textbf{Host:} Copy initial condition to device via
\texttt{cudaMemcpy}\;
Set $\texttt{tpb}=(16,16)$;
$\texttt{bpg}=\bigl(\lceil N_x/16\rceil,\lceil N_y/16
\rceil\bigr)$\;
\For{$t = 0$ \KwTo $T_{\text{end}}$}{
    \textbf{K1:} \texttt{compute\_k\_grid}
    $\langle\!\langle\!\langle\texttt{bpg},\texttt{tpb}
    \rangle\!\rangle\!\rangle$
    --- assign $k_{\text{act}}$ or $k_{\text{pass}}$
    cell-wise\;
    \textbf{K2:} \texttt{compute\_wave\_speeds}
    $\langle\!\langle\!\langle\texttt{bpg},\texttt{tpb}
    \rangle\!\rangle\!\rangle$
    --- compute $\max(|u|+c,\;|v|+c)$ per cell\;
    \textbf{Reduce:} \texttt{cp.max(wave\_speeds)}
    $\Rightarrow a_{\max}$; compute $\Delta t$\;
    \textbf{K3:} \texttt{compute\_fluxes}
    $\langle\!\langle\!\langle\texttt{bpg},\texttt{tpb}
    \rangle\!\rangle\!\rangle$
    --- Rusanov fluxes at all $(N_x{+}1){\times}N_y$
    and $N_x{\times}(N_y{+}1)$ interfaces\;
    \textbf{K4:} \texttt{update\_state}
    $\langle\!\langle\!\langle\texttt{bpg},\texttt{tpb}
    \rangle\!\rangle\!\rangle$
    --- conservative update + dry-bed enforcement\;
}
\textbf{Host:} Copy snapshots from device to host at
checkpoints\;
\end{algorithm}

\subsubsection{Thread Block Configuration}

The two-dimensional domain is tiled with CUDA thread blocks
of dimension $16\times16=256$ threads.
Each thread computes its global cell indices as
\begin{align}
    i &= \text{blockIdx.x}\times\text{blockDim.x}
         + \text{threadIdx.x},
    \label{eq:thread_i} \\
    j &= \text{blockIdx.y}\times\text{blockDim.y}
         + \text{threadIdx.y},
    \label{eq:thread_j}
\end{align}
and a single bounds check prevents out-of-range writes.
For the production grid ($N_x\times N_y=800\times260$) this
yields $50\times17=850$ concurrent thread blocks, sufficient
to saturate the streaming multiprocessor (SM) scheduler of
both the GTX~1650 and the A5000.

\subsection{Shared-Memory Tiling in the Flux Kernel}

The \texttt{compute\_fluxes\_kernel} is the most
computationally intensive stage of each time step.
In the baseline global-memory implementation
(\texttt{gpu\_global.py}), each thread reads six
global-memory values per interface directly from GDDR
memory.
Because the state of cell $(i+1,j)$ is required by both
thread $i$ and thread $i+1$, these loads are redundant
and generate repeated global memory traffic.

The shared-memory variant
(\texttt{gpu\_shared.py}) upgrades the flux kernel with
on-chip tiling.
Each $16\times16$ thread block cooperatively loads a
$(17\times17)$ tile of the $h$, $hu$, $hv$, and
$k_{\text{grid}}$ arrays into \texttt{\_\_shared\_\_}
memory via four distinct passes:
\begin{itemize}[nosep]
    \item All 256 threads load the core $16\times16$
          sub-array in a single coalesced pass.
    \item The rightmost thread column loads the right-halo
          column (required for $x$-fluxes) when
          $i+1<N_x$.
    \item The topmost thread row loads the top-halo row
          (required for $y$-fluxes) when $j+1<N_y$.
    \item A \texttt{\_\_syncthreads()} barrier ensures all
          shared-memory entries are visible before any
          arithmetic proceeds.
\end{itemize}
All subsequent flux arithmetic reads exclusively from the
on-chip tile, aiming to reduce global memory traffic for
this kernel.
The shared-memory variant requires
$4\times17\times17\times\texttt{sizeof(real)}$
bytes of shared memory per block ($\approx4.6$\,KB for
\texttt{float64}), comfortably within the 100\,KB
configurable L1/shared budget per SM on both GPUs.

% ============================================================
\section{Compute-versus-Memory Boundedness Analysis}
\label{sec:roofline}
% ============================================================

\subsection{Roofline Model}

The performance of any GPU kernel is bounded by two
independent hardware ceilings: the floating-point arithmetic
throughput (FLOP/s) and the memory bandwidth (bytes/s).
The roofline model~\cite{Williams2009} characterises a
kernel by its \emph{arithmetic intensity}
\begin{equation}
    I = \frac{\text{FLOP per cell update}}
             {\text{bytes transferred per cell update}},
    \label{eq:arith_intensity}
\end{equation}
and compares it against the \emph{ridge point}
$I^* = \text{FLOP/s}_{\text{peak}}/
\text{BW}_{\text{peak}}$.
Kernels with $I<I^*$ are \emph{memory bound}; kernels
with $I>I^*$ are \emph{compute bound}.

\subsection{Arithmetic Intensity of the Rusanov Kernels}

For the 2D Rusanov stencil in single precision, the state
vector $\bm{U}=(h,hu,hv)^\top$ occupies
$3\times4=12$\,bytes per cell.
Including the $k_{\text{grid}}$ and flux buffer reads and
writes, the total data traffic per cell per time step is
approximately $9\times4=36$\,bytes in single precision
($72$\,bytes in double precision).

The per-cell arithmetic cost is estimated by counting the
operations in \texttt{compute\_fluxes\_kernel} and
\texttt{update\_state\_kernel}: each of the two
contributing interfaces requires approximately 15\,FLOP
(velocity reconstruction, two square roots for wave speeds,
flux evaluations, Rusanov dissipation), and the state update
and source-term evaluation add a further $\approx18$\,FLOP,
giving
\begin{equation}
    I \approx \frac{48\;\text{FLOP}}{36\;\text{bytes}}
            \approx 1.33\;\text{FLOP/byte}.
    \label{eq:intensity_val}
\end{equation}
The GTX~1650 (Turing, FP32) delivers a peak throughput of
$2.98$\,TFLOP/s at a memory bandwidth of $128$\,GB/s,
placing the ridge point at $I^*\approx23.3$\,FLOP/byte.
Since $I\ll I^*$, all physics kernels are strongly
\emph{memory bound} on this architecture.

\subsection{Nsight Compute Profiling}

Quantitative confirmation is provided by NVIDIA Nsight
Compute profiling on the GTX~1650.
Table~\ref{tab:nsight} summarises the average kernel
duration together with hardware compute and memory
throughput utilisation, averaged over repeated invocations
of each kernel during the $800\times260$ production run.

\begin{table}[htbp]
\caption{GTX~1650 Nsight Compute profiling (FP32,
$800\times260$ grid, SpeedOfLight sections).
Duration and throughputs are averaged across
three profiled launches per kernel.}
\label{tab:nsight}
\centering
\begin{tabular}{@{}p{3.0cm}rrrr@{}}
\toprule
\textbf{Kernel}
    & \shortstack[r]{\textbf{Dur.}\\($\mu$s)}
    & \shortstack[r]{\textbf{Comp.}\\(\%)}
    & \shortstack[r]{\textbf{Mem.}\\(\%)}
    & \shortstack[r]{\textbf{BW}\\(GB/s)} \\
\midrule
\texttt{k\_grid}      & 18.8 & 22.4 & 53.9 & 80.2 \\
\texttt{wave\_speeds} & 17.7 & 22.6 & 50.0 & 76.7 \\
\texttt{fluxes}       & 78.8 & 38.1 & 51.9 & 82.7 \\
\texttt{update\_state}& 122.1 & 40.2 & 41.5 & 65.7 \\
\bottomrule
\end{tabular}
\\[3pt]
{\footnotesize
\textit{Abbreviations:}
\texttt{k\_grid} = \texttt{compute\_k\_grid\_kernel};
\texttt{wave\_speeds} = \texttt{compute\_wave\_speeds\_kernel};
\texttt{fluxes} = \texttt{compute\_fluxes\_global\_kernel}.}
\end{table}

The profiling data in Table~\ref{tab:nsight} confirm the
theoretical prediction of Eq.~\eqref{eq:intensity_val}.
Notably, \texttt{update\_state} is the longest-running
kernel (122.1\,$\mu$s) because it reads and writes all
state variables plus all six flux buffer arrays in a single
pass, despite having the lowest memory-throughput
utilisation (41.5\%).
The \texttt{compute\_fluxes\_global} kernel achieves the
highest measured effective bandwidth (82.7\,GB/s,
$\approx65\%$ of the GTX~1650 peak) together with the
highest compute utilisation (38\%), reflecting the
arithmetic-intensive Rusanov dissipation terms.
Across all four kernels, compute utilisation is
consistently below memory utilisation, directly
consistent with the low arithmetic intensity
$I\approx0.3$--$0.7$\,FLOP/byte.
These measurements identify memory bandwidth as the
principal performance-limiting resource and motivate the
SoA layout, single-precision arithmetic, and coalesced
access strategies described in Section~\ref{sec:datalayout}.

% ============================================================
\section{Scaling Analysis}
\label{sec:scaling}
% ============================================================

\subsection{Definitions}

Two complementary scaling studies characterise the parallel
efficiency of the CUDA solver.

\subsubsection{Strong Scaling}

The total problem size $N=N_x\times N_y$ is held fixed
while the number of GPU threads (proportional to $N$)
is implicitly varied by changing the grid resolution.
Strong scaling efficiency is measured as the speedup
relative to the optimised single-threaded C++ reference:
\begin{equation}
    E_{\text{strong}}(N)
    = \frac{T_{\text{serial}}(N)}{T_{\text{GPU}}(N)}.
    \label{eq:strong_eff}
\end{equation}
Perfect strong scaling implies $E_{\text{strong}}$
proportional to $N$.

\subsubsection{Weak Scaling}

Both problem size and GPU thread count increase
proportionally, keeping work per thread constant.
The per-step throughput in million cell updates per second
(MCUPS) is used as the efficiency metric:
\begin{equation}
    \text{MCUPS}
    = \frac{N_x \times N_y \times N_{\text{steps,window}}}
           {T_{\text{window}} \times 10^6}.
    \label{eq:mcups}
\end{equation}
Perfect weak scaling implies constant MCUPS as $N$
increases.
Note that because the CFL condition enforces
$\Delta t\propto\Delta x$, increasing the grid resolution
simultaneously increases the number of time steps required
to reach a fixed physical time, making this a
\emph{simultaneous} spatial and temporal refinement study.

\subsection{Strong Scaling Results}

At small grids ($N_x<200$--$250$ cells, fewer than
$\sim52{,}000$ cells in 2D), kernel launch overhead and
the two-stage global reduction for $\Delta t$ collectively
exceed the useful computation time, so the GPU is slower
than the Numba CPU reference.
As $N_x$ grows beyond this crossover, the GPU's massive
thread-level parallelism saturates the device and speedup
rises steeply.
At the production resolution ($N_x=800$, $208{,}000$
cells), both GPUs deliver approximately $12.4\times$
speedup (GTX~1650: 12.3$\times$; A5000: 12.4$\times$).
The similar speedup ratios reflect that the same
kernel-launch overhead bounds both devices at this
resolution; the A5000 advantage over the GTX~1650 is
captured by the absolute MCUPS metric below.
Dual-panel scaling curves for both GPUs are shown in
Fig.~\ref{fig:dual_scaling} (top row, Section~\ref{sec:multigpu}).

\subsection{Weak Scaling Results}

Grid resolution is increased from $N_x=200$ to
$N_x=3200$ (GTX~1650) and $N_x=6400$ (A5000) in
powers of two while the physical domain
$L_x\times L_y=4.0\times1.3$\,m$^2$ is fixed.
MCUPS is logged every 50 time steps inside
\texttt{run\_gpu\_global}.

MCUPS rises steeply from small grids as thread occupancy
increases, reaching a plateau near the production
resolution ($800\times260$) on the GTX~1650 at
$\approx991$\,MCUPS, and at $\approx6{,}621$\,MCUPS on
the A5000.
The plateau confirms full device occupancy: additional
cells do not impose proportionally higher latency, the
hallmark of good weak scaling on a memory-bound workload.
A modest decline at the largest GTX~1650 grids
($N_x>1200$) is attributed to L2 cache pressure as the
working set
($9\times\texttt{sizeof(real)}\times N_x\times N_y$)
exceeds the 1\,MB on-chip L2.
The A5000 sustains throughput to $N_x=6400$ before any
decline, consistent with its larger L2 and higher
memory bandwidth.
Full MCUPS curves for both GPUs are shown in
Fig.~\ref{fig:dual_scaling} (bottom row,
Section~\ref{sec:multigpu}).

\subsection{Discussion}

For geophysical hazard assessment at high spatial fidelity
($N_x\geq600$, $\Delta x\leq6.7$\,mm), the GPU provides
substantial throughput advantages over CPU-based solvers.
Below the crossover, a single-threaded C++ implementation
is preferable because kernel launch and PCIe transfer
overheads dominate.
The weak scaling plateau confirms that the solver can be
deployed at higher resolutions without proportional loss
of efficiency, subject only to the GPU device memory
limit (24\,GB on the A5000).

% ============================================================
\section{Data Layout Optimisations}
\label{sec:datalayout}
% ============================================================

\subsection{Array of Structs versus Struct of Arrays}

Traditional object-oriented scientific codes group all
physical variables of a single grid cell into a
\emph{struct}:
\begin{verbatim}
struct Cell { float h, hu, hv, k; };
Cell grid[Ny][Nx]; // AoS layout
\end{verbatim}
When a CUDA warp of 32 threads accesses the $h$ field of
32 consecutive cells, the AoS layout forces 32 separate
strided memory transactions because consecutive $h$ values
are separated by \texttt{sizeof(Cell)=16}\,bytes.
Effective memory bandwidth is reduced by a factor of
$\texttt{sizeof(Cell)}/\texttt{sizeof(float)}=4$.

The present solver adopts a \emph{Struct of Arrays} (SoA)
layout:
\begin{verbatim}
float h[Ny*Nx], hu[Ny*Nx];
float hv[Ny*Nx], k_grid[Ny*Nx]; // SoA
\end{verbatim}
When a warp accesses \texttt{h[j*nx+i]} through
\texttt{h[j*nx+i+31]}, all 32 reads target consecutive
4-byte addresses, forming a single 128-byte coalesced
transaction that saturates the GDDR6 memory controller.
In the AoS layout the same 32 reads are separated by
\texttt{sizeof(Cell)=16}\,bytes, forcing 32 strided
transactions and reducing effective bandwidth by a factor
of 4.
The measured SoA throughput advantage at the $3200\times1040$
grid is $1.19\times$ on the GTX~1650 and $1.25\times$ on
the A5000; dual-panel curves appear in
Fig.~\ref{fig:dual_layout} (Section~\ref{sec:multigpu}).

\subsection{Flux Buffer Sizing}

In the Rusanov scheme, fluxes reside on cell
\emph{interfaces}, not cell centres.
For an $N_x\times N_y$ cell domain, the $x$-direction flux
arrays $\hat{\bm{F}}^x$ have size $N_y\times(N_x+1)$
and the $y$-direction arrays $\hat{\bm{G}}^y$ have size
$(N_y+1)\times N_x$.
Allocating each buffer at its natural interface size
eliminates all index-offset arithmetic inside
\texttt{update\_state\_kernel}, reducing per-thread
instruction count and lowering register pressure (the
kernel uses 29 registers per thread, as confirmed by the
Nsight report).

\subsection{Single- versus Double-Precision Strategy}

Consumer GPUs such as the GTX~1650 provide a
$1{:}64$ FP32-to-FP64 throughput ratio.
For the wave-speed-limited Rusanov scheme, the critical
arithmetic is $c=\sqrt{k\,g_z\,h}$, which involves
flow depths of $\mathcal{O}(10^{-2})$\,m.
Round-off errors at FP32 precision are
$\mathcal{O}(10^{-7})$\,m, far below the dry-bed
threshold $h_{\text{dry}}=10^{-4}$\,m, so 32-bit
arithmetic is physically adequate.
\texttt{gpu\_global.py} therefore defaults to
\texttt{float32}, halving the memory footprint per cell
relative to \texttt{float64} and doubling the effective
number of bytes that fit in the L1/L2 caches per SM.
Both GPUs are evaluated in FP32 for benchmark comparisons;
\texttt{float64} is available for high-precision validation
against TITAN2D where truncation artefacts in RMSE metrics
must be minimised.
Table~\ref{tab:precision} quantifies the consequences of
each choice.

\begin{table}[htbp]
\caption{Precision comparison for the GTX~1650 and
A5000 GPU implementations (FP32 vs.\ FP64,
$800\times260$ grid).}
\label{tab:precision}
\centering
\begin{tabular}{@{}p{3.2cm}rr@{}}
\toprule
\textbf{Property}
    & \shortstack[r]{\textbf{GTX~1650}\\\textbf{(FP32)}}
    & \shortstack[r]{\textbf{GTX~1650}\\\textbf{(FP64)}} \\
\midrule
Bytes / cell              & 12    & 24    \\
Peak TFLOP/s              & 2.98  & 0.047 \\
Peak BW (GB/s)            & 128   & 128   \\
MCUPS (800$\times$260)    & 611   & 145   \\
FP32/FP64 speedup         & \multicolumn{2}{c}{$4.2\times$} \\
Mass error (\%)           & 0.0   & 0.0   \\
\bottomrule
\end{tabular}
\end{table}

\subsection{Host--Device Transfer Strategy}

\texttt{cudaMemcpy} transfers are confined to the initial
condition upload at $t=0$ and a checkpoint extraction every
50 time steps (for mass conservation logging and MCUPS
analytics).
No PCIe traffic occurs during the bulk of the time-stepping
loop.
The global reduction \texttt{cp.max(wave\_speeds)} uses a
two-stage CuPy device-resident reduction; only a single
scalar is transferred to the host per step to compute
$\Delta t$, costing approximately $5\,\mu$s --- less than
$5\%$ of the total per-step wall time at the production
grid resolution.

\subsection{Summary of Optimisation Impact}

The Nsight profiling results confirm that all four physics
kernels are memory-bound with compute utilisation
roughly half the memory utilisation (see
Table~\ref{tab:nsight}).
The most productive remaining optimisation pathway is
vectorised loads (\texttt{float4} reads), which would
halve the number of memory transactions per warp and
could raise effective memory bandwidth toward the
hardware peak of $128$\,GB/s on the GTX~1650.

% ============================================================
\section{Extension to Two-Dimensional Flow}
% ============================================================

The one-dimensional framework captures the essential physics of
depth-averaged granular flow but restricts spreading to a single
downslope direction.
Real avalanche events spread laterally as well, driven by internal
pressure gradients in both the downslope ($x$) and cross-slope ($y$)
directions.
This section derives the two-dimensional governing equations,
introduces the dynamic earth-pressure model, and describes the
finite-volume Rusanov scheme used for their numerical solution.

\subsection{Two-Dimensional Governing Equations}

\subsubsection{Coordinate System}

A bed-fitted orthogonal coordinate system is adopted: $x$ is the
downslope direction, $y$ the cross-slope direction, and $z$ normal to
the inclined plane.
The gravitational components resolved along these axes are
\begin{equation}
    g_x = g\sin\zeta, \qquad g_y = 0, \qquad g_z = g\cos\zeta.
    \label{eq:gravity_components}
\end{equation}

\subsubsection{Conservation of Mass}

Depth-averaging the three-dimensional incompressibility condition over
the granular layer yields the two-dimensional continuity equation:
\begin{equation}
    \frac{\partial h}{\partial t}
    + \frac{\partial(hu)}{\partial x}
    + \frac{\partial(hv)}{\partial y} = 0,
    \label{eq:mass_2d}
\end{equation}
where $h(x,y,t)$ is the local flow thickness and $u(x,y,t)$,
$v(x,y,t)$ are the depth-averaged velocity components in the $x$- and
$y$-directions, respectively.

\subsubsection{Conservation of Momentum}

The depth-averaged momentum equations in conservative form read:
\begin{align}
    \frac{\partial(hu)}{\partial t}
    + \frac{\partial}{\partial x}\!\left(hu^2 + \tfrac{1}{2}kg_zh^2\right)
    + \frac{\partial(huv)}{\partial y}
    &= S_x,
    \label{eq:mom_x_2d} \\[4pt]
    \frac{\partial(hv)}{\partial t}
    + \frac{\partial(huv)}{\partial x}
    + \frac{\partial}{\partial y}\!\left(hv^2 + \tfrac{1}{2}kg_zh^2\right)
    &= S_y,
    \label{eq:mom_y_2d}
\end{align}
where $\tfrac{1}{2}kg_zh^2$ represents depth-averaged normal stress
contributions scaled by the earth-pressure coefficient $k$.

\subsubsection{Source Terms}

The source terms $S_x$ and $S_y$ combine gravitational forcing, basal
Coulomb friction, and internal cross-stress:
\begin{equation}
\begin{aligned}
    S_x &= g_xh - \frac{u}{|\bm{u}|}\,hg_z\mu \\
        &\quad - hk\operatorname{sgn}\!\left(\frac{\partial u}
          {\partial y}\right)\sin\phi\,
          \frac{\partial(g_zh)}{\partial y},
\end{aligned}
    \label{eq:source_x}
\end{equation}
\begin{equation}
\begin{aligned}
    S_y &= g_yh - \frac{v}{|\bm{u}|}\,hg_z\mu \\
        &\quad - hk\operatorname{sgn}\!\left(\frac{\partial v}
          {\partial x}\right)\sin\phi\,
          \frac{\partial(g_zh)}{\partial x},
\end{aligned}
    \label{eq:source_y}
\end{equation}
where $|\bm{u}|=\sqrt{u^2+v^2}$ is the local flow speed (regularized
by $\epsilon_r=10^{-3}$ to prevent division by zero) and
$\mu=\tan\delta$ is the basal friction coefficient.
The first term represents gravitational body forcing, the second models
Coulomb friction opposing the local velocity, and the third couples
the two momentum equations through the Mohr--Coulomb internal
cross-stress.

\subsection{Dynamic Earth-Pressure Coefficients}

In two-dimensional flow, different parts of the granular body may
simultaneously be in an active (extending) or passive (compressing)
stress state.
The two states admit distinct coefficients derived from the
Mohr--Coulomb yield criterion:
\begin{equation}
    k = \frac{2}{\cos^2\phi}
        \left(1 \mp \sqrt{1 - \cos^2\phi\left(1+\tan^2\delta\right)}
        \,\right) - 1,
    \label{eq:k_active_passive}
\end{equation}
where the minus sign gives $k_{\text{act}}$ (active, material
extending) and the plus sign gives $k_{\text{pass}}$ (passive,
material compressing).
For $\phi=37.3^\circ$ and $\delta=27.0^\circ$, one obtains
$k_{\text{act}}<1<k_{\text{pass}}$.

\subsubsection{State Selection via Velocity Divergence}

The local stress state is identified at each grid cell from the sign
of the depth-averaged velocity divergence:
\begin{equation}
    \nabla\cdot\bar{\bm{u}}
    = \frac{\partial u}{\partial x} + \frac{\partial v}{\partial y}.
    \label{eq:divergence}
\end{equation}
Expanding flow ($\nabla\cdot\bar{\bm{u}}\geq0$) selects
$k=k_{\text{act}}$; compressing flow selects $k=k_{\text{pass}}$.
Spatial derivatives are computed using central differences at interior
nodes and zero-gradient conditions at boundaries.
The field $k(x,y,t)$ is updated every time step so the pressure
response evolves dynamically.

\subsection{Finite-Volume Rusanov Scheme}

The governing system~\eqref{eq:mass_2d}--\eqref{eq:mom_y_2d} is
written in compact conservative form
\begin{equation}
    \frac{\partial\bm{U}}{\partial t}
    + \frac{\partial\bm{F}(\bm{U})}{\partial x}
    + \frac{\partial\bm{G}(\bm{U})}{\partial y}
    = \bm{S}(\bm{U}),
    \label{eq:conservation_form}
\end{equation}
with state vector $\bm{U}=(h,\,hu,\,hv)^\top$ and flux vectors
\begin{equation}
    \bm{F} = \begin{pmatrix}hu\\ hu^2+\tfrac{1}{2}kg_zh^2\\ huv
             \end{pmatrix}, \quad
    \bm{G} = \begin{pmatrix}hv\\ huv\\ hv^2+\tfrac{1}{2}kg_zh^2
             \end{pmatrix}.
    \label{eq:flux_vectors}
\end{equation}

\subsubsection{Rusanov (Local Lax--Friedrichs) Interface Flux}

At each cell interface, the Rusanov numerical flux is
\begin{equation}
    \hat{\bm{F}}_{i+\nicefrac{1}{2},j}
    = \tfrac{1}{2}\bigl(\bm{F}_L+\bm{F}_R\bigr)
    - \tfrac{1}{2}\,a_{i+\nicefrac{1}{2},j}
      \bigl(\bm{U}_R-\bm{U}_L\bigr),
    \label{eq:rusanov_flux}
\end{equation}
where the local maximum wave speed is
\begin{equation}
    a_{i+\nicefrac{1}{2},j}
    = \max\!\bigl(|u_L|+c_L,\;|u_R|+c_R\bigr),
    \qquad c = \sqrt{kg_zh}.
    \label{eq:rusanov_speed}
\end{equation}
An analogous expression applies to $\hat{\bm{G}}_{i,j+\nicefrac{1}{2}}$
with $v$ replacing $u$.
The Rusanov scheme is first-order accurate and provides strong
stability guarantees without requiring a full Riemann solver, making
it particularly suitable for flows with dry fronts where $h\to0$.

\subsubsection{Dry-Bed Treatment}

Cells with $h<h_{\text{dry}}=10^{-4}$\,m are treated as dry.
Both momentum components are set to zero in such cells, preventing
physically spurious velocity reconstruction near the avalanche margins.

\subsubsection{Time Integration and CFL Condition}

An explicit forward-Euler scheme advances the solution in time:
\begin{equation}
\begin{aligned}
    \bm{U}^{n+1}_{i,j} &= \bm{U}^n_{i,j}
    - \frac{\Delta t}{\Delta x}
      \bigl(\hat{\bm{F}}_{i+\nicefrac{1}{2},j}
           -\hat{\bm{F}}_{i-\nicefrac{1}{2},j}\bigr) \\
    &\quad
    - \frac{\Delta t}{\Delta y}
      \bigl(\hat{\bm{G}}_{i,j+\nicefrac{1}{2}}
           -\hat{\bm{G}}_{i,j-\nicefrac{1}{2}}\bigr)
    + \Delta t\,\bm{S}^n_{i,j}.
\end{aligned}
    \label{eq:fv_update}
\end{equation}
The adaptive time step is determined by the two-dimensional CFL
condition:
\begin{equation}
    \Delta t = \text{CFL} \times \min\!\left(
    \frac{\Delta x}{\max(|u|+c)},\;
    \frac{\Delta y}{\max(|v|+c)}\right),
    \label{eq:cfl_2d}
\end{equation}
with $\text{CFL}=0.3$.

% ============================================================
\section{Two-Dimensional Simulation Setup and TITAN2D Validation}
% ============================================================

\subsection{Physical Parameters and Initial Condition}

The simulated configuration reproduces a hemispherical cap of granular
material released on a planar slope inclined at $\zeta=37.4^\circ$.
Physical and numerical parameters are listed in
Table~\ref{tab:params}.

% FIX: replaced \hline with booktabs rules (\toprule/\midrule/\bottomrule)
%      for publication-standard table formatting.
\begin{table}[htbp]
\caption{Physical and numerical parameters for the 2-D simulation.}
\label{tab:params}
\centering
\begin{tabular}{@{}llcc@{}}
\toprule
\textbf{Parameter}   & \textbf{Symbol}          & \textbf{Value}   & \textbf{Units} \\
\midrule
Slope angle          & $\zeta$                  & $37.4$           & $^\circ$       \\
Basal friction       & $\delta$                 & $27.0$           & $^\circ$       \\
Internal friction    & $\phi$                   & $37.3$           & $^\circ$       \\
Active coefficient   & $k_{\text{act}}$         & computed         & ---            \\
Passive coefficient  & $k_{\text{pass}}$        & computed         & ---            \\
Domain ($x$)         & $L_x$                    & $4.0$            & m              \\
Domain ($y$)         & $L_y$                    & $1.3$            & m              \\
Grid resolution      & $N_x \times N_y$         & $800\times260$   & ---            \\
CFL number           & ---                      & $0.3$            & ---            \\
Dry threshold        & $h_{\text{dry}}$         & $10^{-4}$        & m              \\
\bottomrule
\end{tabular}
\end{table}

The granular pile is initialized as a \emph{cylindrical cap}
(flat disk) of radius $R=0.06$\,m centered at
$(x_c,y_c)=(0.4,0.65)$\,m, with constant height $H_0=R/1.5$:
\begin{equation}
    h(x,y,0) =
    \begin{cases}
        H_0 & (x-x_c)^2+(y-y_c)^2 \leq R^2, \\
        0   & \text{otherwise.}
    \end{cases}
    \label{eq:initial_condition}
\end{equation}
The initial velocity field is zero everywhere.
As time advances, gravitational forcing accelerates the mass
downslope, the pile spreads laterally due to internal pressure
gradients, and Coulomb basal friction progressively decelerates the
flow.

\subsection{Centerline Profile Validation}

TITAN2D is an open-source, parallel finite-volume solver developed
specifically for granular avalanche simulations~\cite{Patra2005}.
It solves the same depth-averaged Mohr--Coulomb system on an adaptive
mesh using a Godunov-type scheme and constitutes a well-established
numerical reference.
Centerline profiles of $h(x,y_c,t)$ are extracted at $y=y_c=0.65$\,m
and compared at five instants: $t=0,\,0.5,\,1.0,\,1.5,\,2.0$\,s
(Fig.~\ref{fig:centerline_val_2d}).

\begin{figure*}[t!]
    \centering
    \includegraphics[width=0.95\textwidth]{centerline_val.pdf}
    \caption{Centerline thickness profiles $h(x,y_c,t)$ at five time
    instants. Dashed lines: TITAN2D reference; solid lines: present
    Dynamic-$k$ Rusanov solver. Colors indicate time: $t=0$\,s
    (black), $t=0.5$\,s (blue), $t=1.0$\,s (green), $t=1.5$\,s
    (orange), $t=2.0$\,s (red).}
    \label{fig:centerline_val_2d}
\end{figure*}
\subsection{Spatial and Temporal Convergence}

A grid convergence and temporal convergence study was conducted to
verify that the reported results are robust with respect to
discretization parameters.
For each study, the centerline RMSE (relative to the TITAN2D
reference at $y_c=0.65$\,m) is computed at $t=1.0$\,s.

For \textbf{spatial convergence}, the grid is refined across four
levels: $100\times32$, $200\times65$, $400\times130$, and
$800\times260$, with the aspect ratio $N_y/N_x\approx0.325$ held
constant and $\text{CFL}=0.3$ fixed.
RMSE decreases monotonically with refinement, confirming that the
first-order Rusanov scheme converges at the expected
$\mathcal{O}(\Delta x)$ rate.
Diminishing returns are observed beyond $N_x=400$, identifying the
$800\times260$ production grid as the optimal balance between
accuracy and computational cost.

For \textbf{temporal convergence}, the grid is fixed at
$400\times130$ and the CFL is varied from $0.8$ down to $0.1$.
Because the first-order Rusanov flux is spatial-error-dominated,
RMSE remains essentially flat across the entire stable CFL envelope.
This confirms that operating at $\text{CFL}=0.3$ delivers full
numerical accuracy without sacrificing GPU throughput.

\begin{figure}[h!]
    \centering
    \includegraphics[width=0.85\textwidth]{grid_convergence.pdf}
    \caption{Spatial (left panel) and temporal (right panel)
    convergence of RMSE against the TITAN2D centerline reference.
    Spatial convergence follows the expected $\mathcal{O}(\Delta x)$
    rate; temporal error is negligible across the stable CFL range
    $0.1$--$0.8$.}
    \label{fig:grid_convergence}
\end{figure}

\subsection{Spatial Footprint and Error Metrics (IoU \& RMSE)}

The numerical footprint is compared against TITAN2D using a wet-cell
threshold $h_{\text{thresh}}=10^{-3}$\,m.
Geometric alignment is evaluated using the Intersection-over-Union
(IoU) metric:
\begin{equation}
    \text{IoU}(t) = \frac{\text{Area}_{\text{Rusanov}}
                    \cap \text{Area}_{\text{TITAN2D}}}
                   {\text{Area}_{\text{Rusanov}}
                    \cup \text{Area}_{\text{TITAN2D}}},
\end{equation}
and the Root-Mean-Square Error (RMSE) of flow depth is computed
exclusively over wetted regions.
Figure~\ref{fig:iou_rmse} illustrates the temporal evolution of both
metrics, demonstrating strong footprint fidelity and bounded depth
error throughout the primary acceleration phase.
The cumulative wetted footprint is shown in
Fig.~\ref{fig:wetted_footprint}.

\begin{figure*}[t!]
    \centering
    \includegraphics[width=0.95\textwidth]{iou_rmse_combined.pdf}
    \caption{Quantitative spatial validation. (Left)
    Intersection-over-Union (IoU) measuring the overlap of the wetted
    footprint between the Dynamic-$k$ Rusanov solver and TITAN2D.
    (Right) RMSE of flow depth $h$ computed on strictly wet cells.}
    \label{fig:iou_rmse}
\end{figure*}

\begin{figure*}[t!]
    \centering
    \includegraphics[width=0.95\textwidth]{wetted_area_comparison.pdf}
    \caption{Cumulative wetted footprint map comparing the maximum
    lateral and longitudinal spread of the Dynamic-$k$ Rusanov solver
    (blue) against TITAN2D (red dashed).}
    \label{fig:wetted_footprint}
\end{figure*}

\subsection{Macroscopic Flow Dynamics ($h_{\max}$ and $\bar{V}$)}

The maximum peak height $h_{\max}$ is tracked continuously to measure
the vertical collapse of the initial cylindrical cap.
The kinetic energy evolution is quantified via a mass-weighted average
velocity:
\begin{equation}
    \bar{V}(t) = \frac{\displaystyle\sum_{i,j} h_{i,j}
                       \sqrt{u_{i,j}^2+v_{i,j}^2}}
                      {\displaystyle\sum_{i,j} h_{i,j}},
\end{equation}
which captures the bulk acceleration of the avalanche under
gravitational forcing, followed by frictional deceleration.
Results are shown in Fig.~\ref{fig:hmax_vel}.

\begin{figure*}[t!]
    \centering
    \includegraphics[width=0.95\textwidth]{hmax_velocity_combined.pdf}
    \caption{Macroscopic flow dynamics. (Left) Collapse of maximum
    flow thickness $h_{\max}(t)$ over time. (Right) Mass-weighted
    average velocity $\bar{V}(t)$ showing the gravitational
    acceleration phase followed by frictional deceleration.}
    \label{fig:hmax_vel}
\end{figure*}

\subsection{Lagrangian Kinematics (Centroid Tracking)}

The trajectory of the avalanche center-of-mass is calculated from the
depth distribution:
\begin{equation}
    X_c(t) = \frac{\displaystyle\sum_{i,j} x_i h_{i,j}}
                  {\displaystyle\sum_{i,j} h_{i,j}}, \qquad
    Y_c(t) = \frac{\displaystyle\sum_{i,j} y_j h_{i,j}}
                  {\displaystyle\sum_{i,j} h_{i,j}}.
\end{equation}
Mapping $\bar{V}$ against the downslope centroid displacement
$d(t)=\sqrt{(X_c-X_{c,0})^2+(Y_c-Y_{c,0})^2}$ provides insight into
run-out distances and momentum conservation
(Fig.~\ref{fig:vel_dist}).

% FIX: was width=0.65\textwidth in a single-column figure environment
%      — overflowed into the adjacent column.  Changed to \columnwidth.
\begin{figure}[htbp]
    \centering
    \includegraphics[width=\columnwidth]{vel_vs_dist_comparison.pdf}
    \caption{Lagrangian trajectory analysis: mass-weighted velocity
    $\bar{V}$ as a function of centroid displacement $d(t)$.
    Selected time stamps are annotated.}
    \label{fig:vel_dist}
\end{figure}

\subsection{Discussion of Results}

The validation suite demonstrates strong agreement between the present
Dynamic-$k$ Rusanov solver and TITAN2D across multiple physical
metrics.

\paragraph{Wave Speed and Profile Shape.}
Both solvers predict the downslope advance of the avalanche front at
nearly identical locations.
The parabolic thinning and progressive flattening of the profile are
captured correctly.
At early times ($t\leq0.5$\,s), the Rusanov solution exhibits minor
numerical diffusion characteristic of first-order fluxes, which
diminishes as the flow transitions to a spreading-dominated regime.

\paragraph{Dynamic $k$ Switching.}
Near the leading edge, material actively extends
($\nabla\cdot\bar{\bm{u}}>0$), selecting $k_{\text{act}}$.
Behind the front, decelerating material compresses, triggering the
local switch to $k_{\text{pass}}$.
This localized stiffening of the pressure wave correctly steepens the
trailing profile slope, validating the 2D divergence field
implementation.

\paragraph{Global Mass Conservation.}
Figure~\ref{fig:mass_2d} confirms that the total volume
$\mathcal{V}(t)=\iint h\,dx\,dy$ remains strictly constant throughout
the simulation, verifying that the discrete finite-volume update
conserves mass to numerical precision.

% FIX: was width=0.7\textwidth in a single-column figure environment.
%      Changed to \columnwidth.
\begin{figure}[htbp]
    \centering
    \includegraphics[width=\columnwidth]{mass_conservation.pdf}
    \caption{Total granular volume $\mathcal{V}(t)$ as a function of
    time, confirming exact mass conservation in the Dynamic-$k$
    Rusanov simulation.}
    \label{fig:mass_2d}
\end{figure}

\subsection{Comparison with the One-Dimensional Model}

Comparing the two-dimensional results with the one-dimensional solver
reveals several important differences.
First, lateral spreading is absent in the 1-D model by construction,
whereas the 2-D simulation shows the pile widening in the
$y$-direction at a rate determined by $kg_z\,\partial h/\partial y$,
which reduces the peak height more rapidly.
Second, the dynamic switching between $k_{\text{act}}$ and
$k_{\text{pass}}$ has a more pronounced effect in 2-D because the
divergence field $\nabla\cdot\bar{\bm{u}}$ captures compression events
that occur even when the downslope component alone is expanding,
producing a richer spatial structure of the pressure coefficient
field.
Third, the Rusanov flux scheme replaces the combination of upwind
advection and centered pressure gradients of the 1-D finite-difference
solver with a unified, conservative numerical flux, ensuring
consistent treatment of shocks and dry fronts within a single
formulation.

% ============================================================
\section{Multi-GPU Performance Comparison}
\label{sec:multigpu}
% ============================================================

All benchmarks reported above were executed on two NVIDIA GPU
architectures: the GTX~1650 (Turing, consumer-class) and the
A5000 (Ampere, workstation-class).
Table~\ref{tab:gpu_specs} summarises the hardware specifications;
Table~\ref{tab:gpu_benchmark} compares peak solver throughput.

\begin{table}[htbp]
\caption{Hardware specifications: GTX~1650 vs.\ A5000.}
\label{tab:gpu_specs}
\centering
\begin{tabular}{@{}p{2.8cm}rr@{}}
\toprule
\textbf{Spec}
    & \shortstack[r]{\textbf{GTX~1650}\\\textbf{(Turing)}}
    & \shortstack[r]{\textbf{A5000}\\\textbf{(Ampere)}} \\
\midrule
CUDA cores         & 896    & 8{,}192 \\
SMs                & 14     & 64      \\
VRAM               & 4\,GB  & 24\,GB  \\
Peak BW (GB/s)     & 128    & 768     \\
Peak FP32 (TFLOP/s)& 2.98   & 27.8    \\
FP64 ratio         & 1{:}64 & 1{:}64  \\
Shared mem./SM     & 64\,KB & 100\,KB \\
Ridge point (F/B)  & 23.3   & 36.2    \\
\bottomrule
\end{tabular}
\end{table}

\begin{table}[htbp]
\caption{Solver benchmark comparison (FP32, \texttt{gpu\_global}).}
\label{tab:gpu_benchmark}
\centering
\begin{tabular}{@{}p{3.4cm}rr@{}}
\toprule
\textbf{Metric}
    & \shortstack[r]{\textbf{GTX~1650}}
    & \shortstack[r]{\textbf{A5000}} \\
\midrule
MCUPS at $800\times260$    & 674    & 3{,}083 \\
Max grid tested            & 3200$\times$1040 & 6400$\times$2080 \\
Peak MCUPS                 & 991    & 6{,}621 \\
GPU/CPU speedup ($800\times260$) & 12.3$\times$ & 12.4$\times$ \\
FP32/FP64 speedup          & 4.2$\times$ & 5.0$\times$ \\
SoA/AoS gain (3200$\times$1040) & 1.19$\times$ & 1.25$\times$ \\
Shared/Global ratio        & 0.92$\times$ & 0.93$\times$ \\
\bottomrule
\end{tabular}
\end{table}

\subsection{Strong and Weak Scaling}

Figure~\ref{fig:dual_scaling} presents dual-panel scaling results.
Both GPUs deliver near-identical CPU-relative speedup (\textasciitilde12.4$\times$)
at $800\times260$, confirming that kernel-launch overhead is
the binding constraint at this size, not compute or bandwidth.
The A5000 sustains higher absolute throughput and scales to
grids that are $4\times$ larger in cell count than the
4\,GB GTX~1650 allows, reaching a plateau of $\sim6{,}600$\,MCUPS
at the $6400\times2080$ configuration.

\begin{figure*}[t!]
    \centering
    \begin{subfigure}[t]{0.48\textwidth}
        \centering
        \includegraphics[width=\textwidth]{results/gtx1650/strong_scaling.pdf}
        \caption{GTX~1650 strong scaling.}
    \end{subfigure}\hfill
    \begin{subfigure}[t]{0.48\textwidth}
        \centering
        \includegraphics[width=\textwidth]{results/a5000/strong_scaling.pdf}
        \caption{A5000 strong scaling.}
    \end{subfigure}
    \\[4pt]
    \begin{subfigure}[t]{0.48\textwidth}
        \centering
        \includegraphics[width=\textwidth]{results/gtx1650/weak_scaling.pdf}
        \caption{GTX~1650 weak scaling (MCUPS).}
    \end{subfigure}\hfill
    \begin{subfigure}[t]{0.48\textwidth}
        \centering
        \includegraphics[width=\textwidth]{results/a5000/weak_scaling.pdf}
        \caption{A5000 weak scaling (MCUPS).}
    \end{subfigure}
    \caption{Dual-GPU scaling study. Top row: speedup relative to
    the Numba CPU reference. Bottom row: solver throughput (MCUPS)
    as a function of grid size. Both GPUs show a~$\sim12.4\times$
    CPU-relative speedup at $800\times260$; the A5000 sustains
    substantially higher absolute throughput and supports larger
    domains due to its 24\,GB VRAM.}
    \label{fig:dual_scaling}
\end{figure*}

\subsection{Precision Impact}

Figure~\ref{fig:dual_precision} compares FP32 and FP64 throughput
on both GPUs.
Although both devices carry a theoretical 1{:}64 FP64-to-FP32
throughput ratio, the measured solver speedup is only 4.2$\times$
(GTX~1650) and 5.0$\times$ (A5000).
This confirms that the solver is memory-bound: the bottleneck is
DRAM bandwidth, not floating-point units, so the full 64$\times$
arithmetic penalty is not observed in practice.

\begin{figure*}[t!]
    \centering
    \begin{subfigure}[t]{0.48\textwidth}
        \centering
        \includegraphics[width=\textwidth]{results/gtx1650/precision_study.pdf}
        \caption{GTX~1650: FP32 vs.\ FP64.}
    \end{subfigure}\hfill
    \begin{subfigure}[t]{0.48\textwidth}
        \centering
        \includegraphics[width=\textwidth]{results/a5000/precision_study.pdf}
        \caption{A5000: FP32 vs.\ FP64.}
    \end{subfigure}
    \caption{Precision impact on solver throughput. FP32 outperforms
    FP64 by 4.2$\times$ (GTX~1650) and 5.0$\times$ (A5000),
    well below the theoretical 64$\times$ peak-FLOP ratio,
    confirming that DRAM bandwidth is the binding resource.}
    \label{fig:dual_precision}
\end{figure*}

\subsection{Data Layout and Shared Memory}

Figure~\ref{fig:dual_layout} compares SoA vs.\ AoS throughput
and shared-memory vs.\ global-memory flux kernels on both GPUs.
The SoA advantage grows with grid size, reaching 1.19$\times$
(GTX~1650) and 1.25$\times$ (A5000) at 3200$\times$1040 as the
AoS layout increasingly saturates the cache with strided
non-coalesced transactions.

The shared-memory flux kernel is consistently \emph{slower} than
the global-memory baseline on both GPUs (ratio 0.92$\times$ on
the GTX~1650, 0.93$\times$ on the A5000 at large grids).
This is a key negative result: for the Rusanov stencil, arithmetic
intensity is too low ($I\approx0.68$\,FLOP/byte) for the
\texttt{\_\_syncthreads()} barrier overhead and extra
shared-memory load passes to pay off.
Shared-memory tiling is more effective for compute-bound stencils
with AI~$>$~1\,FLOP/byte; this solver does not meet that criterion
on current hardware.

\begin{figure*}[t!]
    \centering
    \begin{subfigure}[t]{0.48\textwidth}
        \centering
        \includegraphics[width=\textwidth]{results/gtx1650/data_layout.pdf}
        \caption{GTX~1650: SoA vs.\ AoS.}
    \end{subfigure}\hfill
    \begin{subfigure}[t]{0.48\textwidth}
        \centering
        \includegraphics[width=\textwidth]{results/a5000/data_layout.pdf}
        \caption{A5000: SoA vs.\ AoS.}
    \end{subfigure}
    \\[4pt]
    \begin{subfigure}[t]{0.48\textwidth}
        \centering
        \includegraphics[width=\textwidth]{results/gtx1650/shared_mem_impact.pdf}
        \caption{GTX~1650: Global vs.\ Shared memory flux kernel.}
    \end{subfigure}\hfill
    \begin{subfigure}[t]{0.48\textwidth}
        \centering
        \includegraphics[width=\textwidth]{results/a5000/shared_mem_impact.pdf}
        \caption{A5000: Global vs.\ Shared memory flux kernel.}
    \end{subfigure}
    \caption{Data layout and shared-memory impact on both GPUs.
    Top: SoA consistently outperforms AoS, with the gap widening
    at higher resolutions due to cache saturation.
    Bottom: Shared-memory tiling is slower than global-memory on
    both architectures, a negative result explained by the
    low arithmetic intensity ($I<1$\,FLOP/byte) of the Rusanov
    stencil.}
    \label{fig:dual_layout}
\end{figure*}

\subsection{Block Size Sensitivity}

Figure~\ref{fig:dual_blocksize} shows that the GTX~1650 exhibits
mild sensitivity to block configuration (737--787\,MCUPS,
a 7\% range) with 16$\times$8 threads optimal, while the
A5000 is essentially insensitive ($\leq1$\% variation across all
configurations from $8\times8$ to $32\times32$) due to its larger
SM count and deeper warp pipeline.

\begin{figure*}[t!]
    \centering
    \begin{subfigure}[t]{0.48\textwidth}
        \centering
        \includegraphics[width=\textwidth]{results/gtx1650/block_size_sweep.pdf}
        \caption{GTX~1650: Block size sensitivity.}
    \end{subfigure}\hfill
    \begin{subfigure}[t]{0.48\textwidth}
        \centering
        \includegraphics[width=\textwidth]{results/a5000/block_size_sweep.pdf}
        \caption{A5000: Block size sensitivity.}
    \end{subfigure}
    \caption{Throughput (MCUPS) as a function of CUDA thread block
    configuration. GTX~1650 shows a 7\% variation favoring
    $16\times8$; the A5000 is effectively flat, indicating
    the device is saturated regardless of block geometry.}
    \label{fig:dual_blocksize}
\end{figure*}

\subsection{Cross-GPU Throughput and Roofline}

Figures~\ref{fig:cross_gpu_mcups} and~\ref{fig:cross_gpu_roofline}
present the combined cross-GPU throughput and roofline analyses.
The A5000 delivers $6.7\times$ the MCUPS of the GTX~1650 at peak,
close to the 6$\times$ ratio of their respective memory bandwidths
(768 vs.\ 128\,GB/s), confirming that bandwidth determines
throughput for this workload.

\begin{figure*}[t!]
    \centering
    \begin{subfigure}[t]{0.48\textwidth}
        \centering
        \includegraphics[width=\textwidth]{results/cross_gpu/cross_gpu_mcups.pdf}
        \caption{Peak MCUPS comparison across GPUs.}
    \end{subfigure}\hfill
    \begin{subfigure}[t]{0.48\textwidth}
        \centering
        \includegraphics[width=\textwidth]{results/cross_gpu/cross_gpu_roofline.pdf}
        \caption{Combined roofline: kernel operating points.}
    \end{subfigure}
    \caption{Cross-GPU performance. (Left) Peak solver throughput
    on GTX~1650 and A5000. The A5000 advantage ($\approx6.7\times$)
    closely tracks the bandwidth ratio ($6.0\times$), confirming
    that this solver scales with memory bandwidth.
    (Right) Roofline showing all kernels deeply memory-bound on
    both architectures.}
    \label{fig:cross_gpu_mcups}
\end{figure*}

\begin{figure*}[t!]
    \centering
    \begin{subfigure}[t]{0.48\textwidth}
        \centering
        \includegraphics[width=\textwidth]{results/cross_gpu/cross_gpu_speedup.pdf}
        \caption{GPU-vs-CPU speedup comparison.}
    \end{subfigure}\hfill
    \begin{subfigure}[t]{0.48\textwidth}
        \centering
        \includegraphics[width=\textwidth]{results/cross_gpu/cross_gpu_bandwidth.pdf}
        \caption{Effective memory bandwidth utilisation.}
    \end{subfigure}
    \caption{Cross-GPU speedup and bandwidth efficiency.
    Both GPUs achieve $\sim$12$\times$ speedup over the CPU
    reference at $800\times260$; bandwidth utilisation peaks
    near 65\% of theoretical maximum on both devices.}
    \label{fig:cross_gpu_roofline}
\end{figure*}


% ============================================================
% Nomenclature is unnumbered per conference convention.
\section*{Nomenclature}
% ============================================================

\noindent
\begin{tabular}{@{}ll}
$h$                              & depth-averaged flow thickness (m)              \\
$u,\,v$                          & depth-averaged velocity, $x$- and $y$-dir.\ (m\,s$^{-1}$) \\
$k$                              & earth-pressure coefficient (--)                 \\
$k_{\text{act}},\,k_{\text{pass}}$ & active and passive earth-pressure coeff.\ (--)  \\
$\phi$                           & internal friction angle (rad)                   \\
$\delta$                         & basal friction angle (rad)                      \\
$\zeta$                          & slope inclination angle (rad)                   \\
$\mu$                            & basal friction coefficient, $\mu=\tan\delta$ (--) \\
$g_x,\,g_y,\,g_z$               & gravity components along $x$, $y$, $z$ (m\,s$^{-2}$) \\
$\bm{F},\,\bm{G}$               & flux vectors in $x$- and $y$-directions         \\
$\bm{S}$                         & source term vector                              \\
$\Delta x,\,\Delta y$            & grid spacing (m)                                \\
$\Delta t$                       & time step (s)                                   \\
$h_{\text{dry}}$                 & dry-bed threshold (m)                           \\
\end{tabular}

% ============================================================
\section{Conclusion}
% ============================================================

This paper developed a high-fidelity numerical solver for granular
avalanche motion modeled by the depth-averaged Savage--Hutter
equations.
Theoretical validation against analytical 1D similarity solutions
confirmed accurate mass preservation and spreading dynamics.
The computational framework was then extended to a complete 2D
finite-volume architecture employing a Rusanov Riemann solver with
dynamic active/passive earth-pressure coefficient switching.

Rigorous quantitative benchmarking against the TITAN2D Godunov-type
solver verified the physical accuracy of the 2D model.
The validation suite utilized IoU spatial footprint analysis,
localized RMSE depth tracking, and Lagrangian centroid transport
evaluation, demonstrating excellent agreement in capturing the rapid
acceleration, geometric collapse, lateral spreading, and ultimate
frictional deceleration of the avalanche mass.

A massively parallel GPU architecture was engineered using CUDA to
circumvent the $\mathcal{O}(N_x\times N_y)$ sequential bottleneck
inherent to high-resolution Eulerian grids.
Transitioning state vectors to a Struct of Arrays (SoA) layout and
mapping explicit interface computations directly to parallel hardware
threads allowed the solver to fully saturate GDDR memory bandwidth.
Performance evaluations on the GTX~1650 and A5000 revealed
$\sim12.4\times$ speedup over the optimised CPU reference at the
production grid ($800\times260$), scaling to peak throughputs of
991\,MCUPS (GTX~1650) and 6{,}621\,MCUPS (A5000).
Nsight Compute profiling confirmed all kernels are memory-bound
($I\approx0.3$--$0.7$\,FLOP/byte $\ll I^*$), and the $\sim5\times$
FP32/FP64 speedup ratio (vs.\ theoretical 64$\times$) provides
direct evidence that DRAM bandwidth, not arithmetic throughput,
is the binding resource.
GPU parallelization is therefore a critical enabling technology for
real-time, high-resolution geophysical hazard assessment.

\end{document}